% فصل پنجم: نتیجه‌گیری و پیشنهادها
\chapter{نتیجه‌گیری و پیشنهادها}
\label{ch:conclusion}

در این فصل، ابتدا خلاصه‌ای جامع از فرایند پژوهش و اقدامات انجام‌شده ارائه می‌گردد. سپس دستاوردها و یافته‌های کلیدی پژوهش تبیین شده و در پایان، پیشنهادهایی هدفمند برای ادامه و گسترش این پژوهش در مطالعات آتی ارائه می‌شود.

\section{خلاصه پژوهش}
\label{sec:summary}

در این پایان‌نامه، چارچوب «شبکه‌های عصبی آگاه از فیزیک» (\lr{PINN}) به‌عنوان رویکردی نوین برای حل مسئله مستقیم معادله دیفرانسیل جزئی غیرخطی برگرز مورد بررسی نظری، پیاده‌سازی الگوریتمی و ارزیابی تجربی قرار گرفت.

در بخش نظری، پس از مرور اصول معادلات دیفرانسیل جزئی، روش‌های عددی کلاسیک (تفاضل محدود، اجزای محدود و روش‌های طیفی) و محدودیت‌های آن‌ها در ابعاد بالا، مبانی یادگیری عمیق و سازوکار مشتق‌گیری خودکار تشریح شد. سپس ساختار ریاضی شبکه باقیمانده، نحوه تشکیل تابع هزینه ترکیبی و فرمول‌بندی مدل زمان‌پیوسته تبیین گردید.

در بخش عملی و محاسباتی، مدل زمان‌پیوسته برای معادله غیرخطی برگرز یک‌بعدی با شرایط اولیه سینوسی و مرزهای همگن با استفاده از زبان برنامه‌نویسی پایتون و کتابخانه جکس پیاده‌سازی شد. به‌منظور ارزیابی دقیق نتایج، یک حل‌گر مرجع طیفی مستقل بر پایه روش فوریه و گام‌زنی نمایی مرتبه چهار (\lr{ETDRK4}) توسعه یافت و صحت عملکرد آن با تبدیل نیمه‌تحلیلی کول-هوپف اثبات گردید. آزمایش‌های عددی حل مستقیم با ۲۰۰ نقطه مرزی-اولیه و ۱۰۰۰۰ نقطه هم‌مکانی با راهبرد بهینه‌سازی ترکیبی دو مرحله‌ای آدام و \lr{L-BFGS} به اجرا درآمد و نتایج به دست آمده با مقاله مرجع رئیسی و همکاران (۲۰۱۹) مورد ارزیابی تطبیقی قرار گرفت.

\section{یافته‌ها و دستاوردهای کلیدی}
\label{sec:findings}

مهم‌ترین یافته‌ها و دستاوردهای این پژوهش را می‌توان در محورهای زیر خلاصه نمود:

\begin{enumerate}
\item \textbf{کارایی داده‌ای چشمگیر در حل مسئله مستقیم:} مدل زمان‌پیوسته توانست با بهره‌گیری از تنها ۲۰۰ نقطه داده اولیه و مرزی، میدان پاسخ پیوسته معادله برگرز را بر روی کل دامنه با بیش از ۱۰۰ هزار نقطه بازسازی نموده و به خطای نسبی \FwdRelErr\ در نرم $L_2$ نسبت به حل مرجع دست یابد. این امر مؤید آن است که افزودن قید باقیمانده فیزیکی به تابع هزینه، وابستگی به کلان‌داده‌های آموزشی را برطرف می‌سازد.

\item \textbf{ردیابی دقیق دینامیک موج شوک:} مقایسه مقاطع زمانی در لحظات $t=0.25, 0.50, 0.75$ نشان داد که شبکه عصبی آگاه از فیزیک رفتار فیزیکی تیز شدن پروفیل سرعت و تشکیل موج شوک را با وفاداری بالایی بازتولید می‌نماید.

\item \textbf{تمرکز مکانی خطا در لایه‌های با گرادیان تند:} توزیع مکانی خطای مطلق نشان داد که بیشینه خطا در لایه بسیار باریک پیرامون جبهه شوک متمرکز است. اعمال راهبرد نمونه‌برداری غیریکنواخت و متمرکز در این ناحیه، عاملی تعیین‌کننده در مهار خطا و بهبود دقت تقریب بود.

\item \textbf{اثربخشی راهبرد بهینه‌سازی ترکیبی:} تلفیق الگوریتم آدام (جهت فرار سریع از کمینه‌های محلی آغازین) و بهینه‌ساز شبه‌نیوتنی \lr{L-BFGS} (جهت همگرایی دقیق و عمیق) نقشی حیاتی در رسیدن به دقت‌های مطلوب مهندسی ایفا نمود.

\item \textbf{بازتولیدپذیری با سخت‌افزار استاندارد:} نتایج نشان داد که با استفاده از کتابخانه جکس و دقت ممیز شناور ۶۴ بیتی، امکان بازتولید کیفی و کمی رفتارهای روش‌های فیزیک‌پایه بر روی پردازنده‌های مرکزی استاندارد دانشگاهی بدون نیاز به تجهیزات پردازش فوق‌سریع میسر است.
\end{enumerate}

به‌طور کلی، شبکه‌های آگاه از فیزیک ابزاری قدرتمند و بدون نیاز به شبکه محاسباتی برای حل معادلات غیرخطی هستند و قابلیت ارائه تقریب پیوسته در سراسر فضا-زمان را بدون محدودیت‌های گسسته‌سازی سنتی فراهم می‌سازند.

\section{پیشنهادها برای پژوهش‌های آینده}
\label{sec:future}

به‌منظور ادامه و تعمیق این مسیر پژوهشی، محورهای زیر پیشنهاد می‌گردد:

\begin{enumerate}
\item \textbf{پیاده‌سازی مدل‌های زمان‌گسسته:} پیاده‌سازی و مقایسه کارایی الگوریتم‌های رانگه-کوتای ضمنی با مراحل بالا ($q > 100$) در مقایسه با مدل زمان‌پیوسته در بازه‌های زمانی طولانی.
\item \textbf{توسعه به ابعاد بالاتر و معادلات پیچیده‌تر:} بررسی کارایی روش در حل معادلات برگرز دو‌بعدی و معادلات ناویه‌استوکس و ارزیابی مقیاس‌پذیری الگوریتم در هندسه‌های غیرمنظم.
\item \textbf{روش‌های وزن‌دهی تطبیقی توابع هزینه:} به‌کارگیری راهبردهای وزن‌دهی پویا و گرادیان‌محور برای تنظیم توازن میان مولفه داده ($\mathrm{MSE}_u$) و مولفه فیزیک ($\mathrm{MSE}_f$) در طول آموزش.
\item \textbf{نمونه‌برداری تطبیقی نقاط هم‌مکانی:} توسعه الگوریتم‌های نمونه‌برداری پویا بر اساس گرادیان یا باقیمانده معادله، جهت متمرکزسازی خودکار نقاط هم‌مکانی در نواحی ناپایداری و شوک.
\item \textbf{به‌کارگیری توابع فعال‌ساز تطبیقی:} بررسی تاثیر توابع فعال‌ساز با پارامترهای شیب متغیر (\lr{Adaptive Activation Functions}) جهت بهبود سرعت تقریب جبهه‌های تیز شوک.
\end{enumerate}
